Our results suggest that this task depends a lot on short local patterns in the text. That helps explain why the CNN is already strong at sequence length 64: even with shorter inputs, it can pick up useful topic phrases and keywords. So the CNN advantage at short length is not only a score difference, but also a good match between the model and the kind of signal in AG News.

At the same time, the BiLSTM improves more as we increase sequence length, and at 256 tokens it almost catches the CNN (Table~\ref{tab:test_results}). A likely reason is that the extra context helps the LSTM in examples where early words are not enough to decide the class. This seems especially relevant for Business and Sci/Tech, where articles often contain mixed vocabulary. Still, even at 256 the LSTM does not clearly pass the CNN, which suggests that local cues remain very important for this dataset.

The class-wise results support this broader picture. World and Sports are generally easier for both models, while Business and Sci/Tech remain harder. This tells us that the remaining errors are mostly about separating semantically close topics, not about basic model failure.

Overall, CNN-seq256 is still the best single configuration in our experiments, but the margin over LSTM-256 is small. Although the CNN is better, it does not scale with sequence length as well as the LSTM, we think that if sequence length were increased further, the LSTM might eventually outperform the CNN. This is because the LSTM can potentially leverage longer-range dependencies and more global context, while the CNN relies more on local patterns that may saturate in usefulness after a certain point. In future work, it would be interesting to explore even longer sequence lengths and see if the LSTM can eventually surpass the CNN as it gets more context to work with.

